% Options for packages loaded elsewhere
\PassOptionsToPackage{unicode}{hyperref}
\PassOptionsToPackage{hyphens}{url}
\documentclass[
  ignorenonframetext,
]{beamer}
\newif\ifbibliography
\usepackage{pgfpages}
\setbeamertemplate{caption}[numbered]
\setbeamertemplate{caption label separator}{: }
\setbeamercolor{caption name}{fg=normal text.fg}
\beamertemplatenavigationsymbolsempty
% remove section numbering
\setbeamertemplate{part page}{
  \centering
  \begin{beamercolorbox}[sep=16pt,center]{part title}
    \usebeamerfont{part title}\insertpart\par
  \end{beamercolorbox}
}
\setbeamertemplate{section page}{
  \centering
  \begin{beamercolorbox}[sep=12pt,center]{section title}
    \usebeamerfont{section title}\insertsection\par
  \end{beamercolorbox}
}
\setbeamertemplate{subsection page}{
  \centering
  \begin{beamercolorbox}[sep=8pt,center]{subsection title}
    \usebeamerfont{subsection title}\insertsubsection\par
  \end{beamercolorbox}
}
% Prevent slide breaks in the middle of a paragraph
\widowpenalties 1 10000
\raggedbottom
\AtBeginPart{
  \frame{\partpage}
}
\AtBeginSection{
  \ifbibliography
  \else
    \frame{\sectionpage}
  \fi
}
\AtBeginSubsection{
  \frame{\subsectionpage}
}
\usepackage{iftex}
\ifPDFTeX
  \usepackage[T1]{fontenc}
  \usepackage[utf8]{inputenc}
  \usepackage{textcomp} % provide euro and other symbols
\else % if luatex or xetex
  \usepackage{unicode-math} % this also loads fontspec
  \defaultfontfeatures{Scale=MatchLowercase}
  \defaultfontfeatures[\rmfamily]{Ligatures=TeX,Scale=1}
\fi
\usepackage{lmodern}
\ifPDFTeX\else
  % xetex/luatex font selection
\fi
% Use upquote if available, for straight quotes in verbatim environments
\IfFileExists{upquote.sty}{\usepackage{upquote}}{}
\IfFileExists{microtype.sty}{% use microtype if available
  \usepackage[]{microtype}
  \UseMicrotypeSet[protrusion]{basicmath} % disable protrusion for tt fonts
}{}
\makeatletter
\@ifundefined{KOMAClassName}{% if non-KOMA class
  \IfFileExists{parskip.sty}{%
    \usepackage{parskip}
  }{% else
    \setlength{\parindent}{0pt}
    \setlength{\parskip}{6pt plus 2pt minus 1pt}}
}{% if KOMA class
  \KOMAoptions{parskip=half}}
\makeatother
\usepackage{graphicx}
\makeatletter
\newsavebox\pandoc@box
\newcommand*\pandocbounded[1]{% scales image to fit in text height/width
  \sbox\pandoc@box{#1}%
  \Gscale@div\@tempa{\textheight}{\dimexpr\ht\pandoc@box+\dp\pandoc@box\relax}%
  \Gscale@div\@tempb{\linewidth}{\wd\pandoc@box}%
  \ifdim\@tempb\p@<\@tempa\p@\let\@tempa\@tempb\fi% select the smaller of both
  \ifdim\@tempa\p@<\p@\scalebox{\@tempa}{\usebox\pandoc@box}%
  \else\usebox{\pandoc@box}%
  \fi%
}
% Set default figure placement to htbp
\def\fps@figure{htbp}
\makeatother
\setlength{\emergencystretch}{3em} % prevent overfull lines
\providecommand{\tightlist}{%
  \setlength{\itemsep}{0pt}\setlength{\parskip}{0pt}}
% Slide layout defaults for warning-clean scientific decks.
\usepackage{etoolbox}
\IfFileExists{xurl.sty}{\usepackage{xurl}}{}
\IfFileExists{seqsplit.sty}{\usepackage{seqsplit}}{\newcommand{\seqsplit}[1]{#1}}
\protected\def\breakseq#1{\seqsplit{#1}}
\protected\def\breaktt#1{\begingroup\ttfamily\seqsplit{#1}\endgroup}
\setlength{\emergencystretch}{6em}
\tolerance=5000
\hbadness=10000
\hfuzz=1pt
\setlength{\tabcolsep}{2pt}
\AtBeginEnvironment{longtable}{\tiny\renewcommand{\arraystretch}{0.86}\setlength{\tabcolsep}{1pt}}
\AtBeginEnvironment{tabular}{\tiny\renewcommand{\arraystretch}{0.86}\setlength{\tabcolsep}{1pt}}

% Natbib and cross-reference fallbacks — slides don't load natbib
% or cleveref, but combined-PDF manuscript prose may emit these
% commands. The fallback renders citations as a bracketed key list
% and cross-references as detokenized labels so slides don't fail on
% undefined control sequences or raw underscores. \providecommand is
% a no-op if packages load later.
\providecommand{\citep}[1]{[#1]}
\providecommand{\citet}[1]{#1}
\providecommand{\citealp}[1]{#1}
\providecommand{\citeauthor}[1]{#1}
\providecommand{\citeyear}[1]{#1}
\providecommand{\cref}[1]{\texttt{\detokenize{#1}}}
\providecommand{\Cref}[1]{\texttt{\detokenize{#1}}}

\newtheorem{proposition}{Proposition}
\newtheorem{hypothesis}{Hypothesis}
\usepackage{bookmark}
\IfFileExists{xurl.sty}{\usepackage{xurl}}{} % add URL line breaks if available
\urlstyle{same}
% fallback for those not using the hyperref driver hyperxmp:
\makeatletter
\@ifundefined{xmpquote}{\newcommand{\xmpquote}[1]{#1}}{}
\makeatother
\hypersetup{
  hidelinks,
  pdfcreator={LaTeX via pandoc}}

\author{\texorpdfstring{}{}}
\date{}

\begin{document}

\section{Data description}\label{sec:data}

\begin{frame}[fragile,allowframebreaks]{Data description}
The demonstration dataset consists of two comma-separated files under
\texttt{data/fixtures/}, described by the machine-readable descriptor
\breaktt{data/example\_descriptor.json}:

\begin{itemize}
\tightlist
\item
  \breaktt{fixtures/measurements.csv} --- one row per synthetic sample,
  with a bounded measurement, an assignment group, a capture instrument,
  and a collection date.
\item
  \breaktt{fixtures/subjects.csv} --- one row per synthetic subject, with
  an enrollment site and date. Each measurement references a subject
  through the \texttt{subject\_id} foreign key.
\end{itemize}
\end{frame}

\begin{frame}[fragile,allowframebreaks]{File inventory}
\protect\phantomsection\label{file-inventory}
The descriptor declares each file's media type, sha256 checksum, and row
count. \breakseq{[@fig:file\_inventory]} shows the declared row counts, read
directly from the descriptor by \breaktt{file\_inventory\_rows()} in
\breaktt{src/data\_descriptor/figures.py} and plotted by the thin
analysis script
\href{../scripts/generate_figures.py}{\breaktt{scripts/generate\_figures.py}}.

\begin{figure}
\centering
\pandocbounded{\includegraphics[keepaspectratio,alt={File inventory: declared row counts per file, from file\_inventory\_rows(). The measurement table is the larger of the two files; the subject table is its dimension companion. Both are declared as text/csv. The bar values are the row counts recorded in the descriptor, not re-derived here --- the quality gate (\breakseq{[@sec:quality]}) is where declared and actual counts are reconciled.}]{../figures/file_inventory.png}}
\caption{File inventory: declared row counts per file, from
\protect\breaktt{file\_inventory\_rows()}. The measurement table is the
larger of the two files; the subject table is its dimension companion.
Both are declared as \protect\texttt{text/csv}. The bar values are the
row counts recorded in the descriptor, not re-derived here --- the
quality gate (\breakseq{[@sec:quality]}) is where declared and actual counts
are reconciled.}\label{fig:file_inventory}
\end{figure}

Running the script regenerates the figures under
\breaktt{manuscript/figures/}; the prose below and in later sections
describes what those artifacts show. Because the figures are produced
from the descriptor and the fixture bytes, they cannot drift from the
data without the tests and the quality gate noticing.
\end{frame}

\begin{frame}[allowframebreaks]{Release boundary}
\protect\phantomsection\label{release-boundary}
The descriptor is deliberately a \emph{metadata} object. For a real
dataset, the descriptor and its checksums can be published even when the
underlying bytes are access-controlled: the checksums let a downstream
user verify integrity once they obtain the files through the appropriate
channel. In this template the fixture bytes are public and small, so
they are shipped alongside the descriptor and used to demonstrate
end-to-end verification.
\end{frame}

\end{document}
